\chapter{1986年上海卷（文）}

\section{单选题}

\begin{problem}

平移坐标轴，把原点 \(O(0,0)\) 移到 \(O'(2,-1)\)，则点 \((-1,-3)\) 在新坐标系中的坐标是（\quad）

\choices
    {\(\left( 3,2 \right)\)}
    {\(\left( -3,-2 \right)\)}
    {\(\left( -3,2 \right)\)}
    {\(\left( 3,-2 \right)\)}

\end{problem}

\begin{answer}
B
\end{answer}

\begin{problem}

下列三角关系式中不正确的是（\quad）

\choices
    {\(\sin\alpha+\sin\beta=2\sin\frac{\beta+\alpha}{2}\cos\frac{\beta-\alpha}{2} \)}
    {\(\sin\alpha-\sin\beta=2\cos\frac{\beta+\alpha}{2}\sin\frac{\beta-\alpha}{2} \)}
    {\(\cos \alpha + \cos \beta = 2\cos \frac{\beta + \alpha}{2}\cos \frac{\beta - \alpha}{2}\)}
    {\(\cos \alpha - \cos \beta = 2\sin \frac{\beta + \alpha}{2}\sin \frac{\beta - \alpha}{2}\)}

\end{problem}

\begin{answer}
B
\end{answer}

\begin{problem}

在直角坐标系中，直线 \(x + \sqrt{3} y + 1 = 0\) 的倾角是（\quad）

\choices
    {\(\frac{1}{6}\pi\)}
    {\(\frac{1}{3}\pi\)}
    {\(\frac{2}{3}\pi\)}
    {\(\frac{5}{6}\pi\)}

\end{problem}

\begin{answer}
D
\end{answer}

\begin{problem}

函数 \(y = 2\tan (3x + \frac{1}{6}\pi) \) 的最小正周期是（\quad）

\choices
    {\(\frac{1}{6}\pi\)}
    {\(\frac{1}{3}\pi\)}
    {\(\frac{1}{2}\pi\)}
    {\(\frac{2}{3}\pi\)}

\end{problem}

\begin{answer}
B
\end{answer}

\begin{problem}

在空间，下列命题中正确的是（\quad）

\choices
    {如果两直线 \(a,b\) 与直线 \(l\) 所成的角相等，那么 \(a\parallel b\)}
    {如果两直线 \(a,b\) 与平面 \(\alpha\) 所成的角相等，那么 \(a\parallel b\)}
    {如果直线 \(l\) 与两平面 \(\alpha,\beta\) 所成的角都是直角，那么 \(\alpha\parallel\beta\)}
    {如果平面 \(\gamma\) 与两平面 \(\alpha,\beta\) 所成的二面角都是直二面角，那么 \(\alpha\parallel\beta\)}

\end{problem}

\begin{answer}
C
\end{answer}

\begin{problem}

函数 \(y = \tan\left(x + \frac{1}{3}\pi\right)\) 的定义域是（\quad）

\choices
    {\(\{x \mid x \in \R \text{ 且 } x \neq k\pi + \frac{1}{6}\pi,\ k \in \Z\}\)}
    {\(\{x \mid x \in \R \text{ 且 } x \neq k\pi - \frac{1}{6}\pi,\ k \in \Z\}\)}
    {\(\{x \mid x \in \R \text{ 且 } x \neq 2k\pi + \frac{1}{6}\pi,\ k \in \Z\}\)}
    {\(\{x \mid x \in \R \text{ 且 } x \neq 2k\pi - \frac{1}{6}\pi,\ k \in \Z\}\)}

\end{problem}

\begin{answer}
A
\end{answer}

\begin{problem}

若空间有10个点，则可以确定的平面总数最多有（\quad）

\choices
    {90个}
    {100个}
    {120个}
    {150个}

\end{problem}

\begin{answer}
C
\end{answer}

\begin{problem}

下列双曲线方程中，以 \(y = \pm x \) 为渐近线的是（\quad）

\choices
    {\(\frac{x^{2}}{16}-\frac{y^{2}}{4}=1 \)}
    {\(\frac{x^{2}}{4}-\frac{y^{2}}{16}=1 \)}
    {\(\frac{x^{2}}{2}-\frac{y^{2}}{1}=1 \)}
    {\(\frac{x^{2}}{1}-\frac{y^{2}}{2}=1 \)}

\end{problem}

\begin{answer}
A
\end{answer}

\begin{problem}

复数 \(\sqrt{3} - \i\) 的幅角的主值是（\quad）

\choices
    {\(\frac{1}{6}\pi\)}
    {\(\frac{5}{6}\pi\)}
    {\(7\pi /6\)}
    {\(11\pi /6\)}

\end{problem}

\begin{answer}
D
\end{answer}

\begin{problem}

等式 \(\log_3 x^2 = 2\) 成立是等式 \(\log_3 x = 1\) 成立的（\quad）

\choices
    {充分条件但不是必要条件}
    {必要条件但不是充分条件}
    {充分必要条件}
    {既不充分又不必要的条件}

\end{problem}

\begin{answer}
B
\end{answer}

\begin{problem}

在等比数列中，已知首项为 \(\frac{9}{8}\)，末项为 \(\frac{1}{3}\)，公比为 \(\frac{2}{3}\)，则项数是（\quad）

\choices
    {3}
    {4}
    {5}
    {6}

\end{problem}

\begin{answer}
B
\end{answer}

\begin{problem}

已知定直线 \(l\) 和外一定点 \(A\)，过点 \(A\) 且与 \(l\) 相切的圆的圆心轨迹是（\quad）

\choices
    {抛物线}
    {双曲线}
    {椭圆}
    {直线}

\end{problem}

\begin{answer}
A
\end{answer}

\begin{problem}

若 \(\sin \frac{1}{2}\theta = \frac{3}{5}\)，\(\cos \frac{1}{2}\theta = -\frac{4}{5}\)，则 \(\theta\) 角的终边在（\quad）

\choices
    {第一象限}
    {第二象限}
    {第三象限}
    {第四象限}

\end{problem}

\begin{answer}
D
\end{answer}

\begin{problem}

下列不等式中正确的是（\quad）

\choices
    {\(\sin\left(-\frac{5}{7}\pi\right)>\sin\left(\frac{4}{7}\pi\right)\)}
    {\(\tan\left(\frac{15\pi}{8}\right)>\tan\left(-\frac{3}{8}\pi\right)\)}
    {\(\sin\left(-\frac{1}{3}\pi\right)>\sin\left(-\frac{1}{6}\pi\right)\)}
    {\(\cos\left(-\frac{8}{5}\pi\right)>\cos\left(-\frac{9}{4}\pi\right)\)}

\end{problem}

\begin{answer}
B
\end{answer}

\begin{problem}

集合 \(\{a, b, c\}\) 的所有子集的个数是（\quad）

\choices
    {5}
    {6}
    {7}
    {8}

\end{problem}

\begin{answer}
D
\end{answer}

\begin{problem}

设自变量 \(x \in \R\)，下列各函数中奇函数是（\quad）

\choices
    {\(y = x + 3 \)}
    {\(y = -2x^{2} \)}
    {\(y = \lg 3^x \)}
    {\(y = -|x| \)}

\end{problem}

\begin{answer}
C
\end{answer}

\begin{problem}

\(a, b\) 满足 \(0 < a < b < 1\). 下列不等式中正确的是（\quad）

\choices
    {\(a^a < a^b\)}
    {\(b^a < b^b\)}
    {\(a^a < b^a\)}
    {\(b^a < a^b\)}

\end{problem}

\begin{answer}
C
\end{answer}

\begin{problem}

正方体 \(ABCD - A_{1}B_{1}C_{1}D_{1}\) 中，直线 \(BC_{1}\) 与 \(AC\)（\quad）

\choices
    {相交且垂直}
    {相交但不垂直}
    {异面且垂直}
    {异面但不垂直}

\end{problem}

\begin{answer}
D
\end{answer}

\begin{problem}

设复数 \(z = a + b\i\)（\(a,b \in \R\)，且 \(b \neq 0\)），则 \(|z^2|\)、\(|z|^2\)、\(z^2\) 的关系是（\quad）

\choices
    {\(|z^{2}|=|z|^{2}\neq z^{2} \)}
    {\(|z^{2}|=|z|^{2}=z^{2} \)}
    {\(|z^{2}| \neq |z|^{2} = z^{2} \)}
    {互不相等}

\end{problem}

\begin{answer}
A
\end{answer}

\begin{problem}

过抛物线焦点 \(F\) 的直线与抛物线相交于 \(A\)、\(B\) 两点，若 \(A\)、\(B\) 在抛物线准线上的射影分别是 \(A_{1}\)、\(B_{1}\)，则 \(\angle A_{1}FB_{1}\) 等于（\quad）

\choices
    {\(45^{\circ}\)}
    {\(60^{\circ}\)}
    {\(90^{\circ}\)}
    {\(120^{\circ}\)}

\end{problem}

\begin{answer}
C
\end{answer}

\section{填空题}

\begin{problem}

若三点 \(P_{1}\)、\(P_{2}\) 和 \(P\) 在一条直线上，点 \(P_{1}\) 和点 \(P_{2}\) 在直角坐标系中的坐标分别为 \((0,-6)\) 和 \((3,0)\)，且 \(\frac{P_{1}P}{PP_{2}}=-\frac{1}{2}\)，则点 \(P\) 的坐标是\fillinblank{}.

\end{problem}

\begin{answer}
\(\left( -3,-12 \right)\)
\end{answer}

\begin{problem}

正方体的全面积是 \(6\)，则其内切球的体积是\fillinblank{}.

\end{problem}

\begin{answer}
\(\frac{\pi}{6}\)
\end{answer}

\begin{problem}

设直线 \(l\) 在 \(x\) 轴和 \(y\) 轴上的截距分别为 \(4\) 和 \(3\)，则点 \((2, -1)\) 到 \(l\) 的距离是\fillinblank{}.

\end{problem}

\begin{answer}
2
\end{answer}

\begin{problem}

\(\displaystyle\lim_{n\to \infty}\frac{1 + 3 + 5 + \cdots + (2n - 1)}{n(n + 1)} =\)\fillinblank{}

\end{problem}

\begin{answer}
1
\end{answer}

\begin{problem}

\(\left( \sqrt[3]{x}-\frac{1}{\sqrt{x}} \right)^8\) 的展开式中，\(x\) 的一次项的系数是\fillinblank{}.

\end{problem}

\begin{answer}
28或\(\mathrm C_8^6\)
\end{answer}

\begin{problem}

不等式 \(\sqrt{x - 1} < 3\) 的解是\fillinblank{}.

\end{problem}

\begin{answer}
\(1\leqslant x<10\)
\end{answer}

\begin{problem}

轴截面是等边三角形的圆锥，它的侧面展开图扇形的圆心角的弧度数等于\fillinblank{}.

\end{problem}

\begin{answer}
\(\pi\)
\end{answer}

\begin{problem}

用 \(1\)、\(2\)、\(3\)、\(4\) 四个数字组成没有重复数字的四位奇数的个数是\fillinblank{}.

\end{problem}

\begin{answer}
\(12\)
\end{answer}

\begin{problem}

已知 \(\cos \theta = -\frac{3}{5}\) 且 \(\pi < \theta < 3\pi / 2\)，则 \(\cot(\theta - \frac{1}{2}\pi) =\)\fillinblank{}.

\end{problem}

\begin{answer}
7
\end{answer}

\begin{problem}

函数 \(y = \log_2(\log_{\frac{1}{2}}x) \) 的定义域是\fillinblank{}.

\end{problem}

\begin{answer}
\((0,1)\)
\end{answer}

\section{解答题}

\begin{problem}

设 \(f(x)=\frac{1}{2}\log_{\frac{1}{3}}(x+1)\)，\(g(x)=\log_{\frac{1}{3}}(1-x)\)，试解不等式 \(f(x)\geqslant g(x)\).

\end{problem}

\begin{solution}
%
由
\[
\frac{1}{2}\log_{\frac{1}{3}}(x+1)\geqslant \log_{\frac{1}{3}}(1-x)
\]
得
\[
\begin{cases}
x+1>0,\\
1-x>0,\\
\log_{\frac{1}{3}}(x+1)\geqslant \log_{\frac{1}{3}}(1-x)^2.
\end{cases}
\]
解不等式\circled{1}，得 \(x>-1\)，
解不等式\circled{2}，得 \(x<1\)，
解不等式\circled{3}，得 \(x+1\leqslant(1-x)^2\)，即
\[
x\leqslant 0 \text{ 或 } x\geqslant 3.
\]
所以原不等式的解是 \(-1<x\leqslant 0\).
\end{solution}

\begin{problem}

设 \(y=\frac{1-\sin\theta\cos\theta}{1+\sin\theta\cos\theta} \) ，当 \(\theta\in[0,\pi] \) 上分别取何值时，\(y\) 取得最大值和最小值.

\end{problem}

\begin{solution}
%
由
\[
y=\frac{1-\sin\theta\cos\theta}{1+\sin\theta\cos\theta}=\frac{2-\sin 2\theta}{2+\sin 2\theta}.
\]
因此当 \(\theta=\frac{3}{4}\pi\) 时，\(\sin 2\theta=-1\)，\(y\) 取得最大值；当 \(\theta=\frac{1}{4}\pi\) 时，\(\sin 2\theta=1\)，\(y\) 取得最小值.
\end{solution}

\begin{problem}

写出棣莫佛定理并用数学归纳法加以证明.

\end{problem}

\begin{solution}
%
棣莫佛定理：复数的 \(n\) 次幂（\(n\) 是正整数）的模等于这个复数的模的 \(n\) 次幂，它的幅角等于这个复数的幅角的 \(n\) 倍。若记 \(z=r(\cos \theta+\i\sin \theta)\)，\(n\) 为正整数，则
\[
z^{n}=r^{n}\left( \cos n\theta+\i\sin n\theta \right).
\]

证明：

（1）当 \(n=1\) 时，显然成立。

（2）假设当 \(n=k\) 时定理成立，即有
\[
z^{k}=r^{k}\left( \cos k\theta+\i\sin k\theta \right).
\]
当 \(n=k+1\) 时，
\[
\begin{aligned}
z^{k+1}
&=z^{k}\cdot z\\
&=r^{k}\left( \cos k\theta+\i\sin k\theta \right)\cdot r\left( \cos \theta+\i\sin \theta \right)\\
&=r^{k+1}\left[ \left( \cos k\theta\cos \theta-\sin k\theta\sin \theta \right)+\i\left( \sin k\theta\cos \theta+\cos k\theta\sin \theta \right) \right]\\
&=r^{k+1}\left[ \cos \left( k+1 \right)\theta+\i\sin \left( k+1 \right)\theta \right].
\end{aligned}
\]
即定理在 \(n=k+1\) 时也成立。根据（1）和（2），对于任何正整数 \(n\)，定理都成立。
\end{solution}

\begin{problem}

在正方体 \(ABCD - A_{1}B_{1}C_{1}D_{1}\) 中， \(E\) 为棱 \(CC_{1}\) 的中点，求截面\(A_{1}BD\) 和截面 \(EBD\) 所成二面角的度数.

\begin{center}
\adjustbox{max width=0.42\linewidth,max totalheight=4.8cm}{%
\includegraphics[page=4,pagebox=mediabox,trim=180pt 267.6pt 295.6pt 392pt,clip]{source/普通高考/1986/1986上海文.pdf}%
}
\end{center}

\end{problem}

\begin{solution}
%
连 \(AC\)，与 \(BD\) 交于 \(O\)，则 \(O\) 为 \(AC\)、\(BD\) 的中点。连 \(A_{1}O\)、\(EO\)、\(A_{1}E\) 和 \(A_{1}C_{1}\).
\[
\text{因为 } A_{1}B=A_{1}D,
\]
\[
EB=ED,
\]
\[
\text{所以 } A_{1}O\perp BD,
\]
且 \(EO\perp BD\).

故 \(\angle A_{1}OE\) 为两截面所成二面角的平面角。设 \(\angle A_{1}OE=\alpha\)，正方体的棱长为 \(a\).
\[
\text{因为 } A_{1}O=\sqrt{AA_{1}^{2}+AO^{2}}=\frac{\sqrt{6}a}{2},
\]
\[
EO=\sqrt{EC^{2}+CO^{2}}=\frac{\sqrt{3}a}{2},
\]
\[
A_{1}E=\sqrt{A_{1}C_{1}^{2}+C_{1}E^{2}}=\frac{3a}{2},
\]
在 \(\triangle A_{1}OE\) 中，
\[
\cos \alpha=\frac{A_{1}O^{2}+EO^{2}-A_{1}E^{2}}{2A_{1}O\cdot EO}.
\]
所以 \(\alpha=90^{\circ}\)，即两截面所成的二面角为 \(90^{\circ}\).
\end{solution}

\begin{problem}

如果抛物线 \(y^{2}=px\) 和圆 \((x-2)^{2}+y^{2}=3\) 相交，它们在 \(x\) 轴上方交的交点为 \(A\) 和 \(B\)，那么当 \(p\) 为何值时，线段 \(AB\) 的中点 \(M\) 在直线 \(y=x\) 上.

\bitmapfigure[max width=0.42\linewidth,max totalheight=4.8cm]{img/1986/shanghai_science/q6_fig1.png}

\end{problem}

\begin{solution}
%
设点 \(A\) 的坐标为 \((x_{1},y_{1})\)，点 \(B\) 的坐标为 \((x_{2},y_{2})\).
\[
y^{2}=px, \circled{1}
\]
\[
(x-2)^{2}+y^{2}=3. \circled{2}
\]
\circled{1} 代入 \circled{2}，得
\[
x^{2}+(p-4)x+1=0.
\]
所以
\[
x_{1}+x_{2}=4-p,\qquad x_{1}x_{2}=1.
\]

设 \(M\) 点坐标为 \((x_{M},y_{M})\)，则
\[
x_{M}=\frac{x_{1}+x_{2}}{2}=\frac{4-p}{2}.
\]
因为点 \(A\) 和点 \(B\) 在 \(x\) 轴上方，
\[
\begin{aligned}
y_{1}+y_{2}
&=\sqrt{\left( y_{1}+y_{2} \right)^{2}}\\
&=\sqrt{y_{1}^{2}+y_{2}^{2}+2y_{1}y_{2}}\\
&=\sqrt{p\left( x_{1}+x_{2} \right)+2p\sqrt{x_{1}x_{2}}}\\
&=\sqrt{p(4-p)+2p}=\sqrt{6p-p^{2}}.
\end{aligned}
\]
于是
\[
y_{M}=\frac{y_{1}+y_{2}}{2}=\frac{\sqrt{6p-p^{2}}}{2}.
\]
因为点 \(M\) 在直线 \(y=x\) 上，
\[
\sqrt{6p-p^{2}}=4-p.
\]
即
\[
6p-p^{2}=\left( 4-p \right)^{2},
\]
所以
\[
p^{2}-7p+8=0.
\]
解得
\[
p=\frac{7\pm \sqrt{17}}{2}.
\]
又因为 \(4-p\geqslant 0\)，故舍去 \(p=\frac{7+\sqrt{17}}{2}\). 因此
\[
p=\frac{7-\sqrt{17}}{2}.
\]
\end{solution}
